\Article{7}
In place of formula (3), which expresses the weight of the most plausible determination, several others can be given that will be worth developing.

First we observe that if equations (2) are multiplied respectively by \(\frac ap,\ \frac{a'}{p'},\ \frac{a''}{p''}\) etc. and added, the result is
\[\relax
[aa]x^0+[ab]y^0+[ac]z^0+\text{etc.}=\frac{aL}{p}+\frac{a'L'}{p'}+\frac{a''L''}{p''}+\text{etc.}
\]
The left side becomes \(=0\); by analogy, we denote the right side by \([aL]\); we therefore have
\[\relax
[aL]=0,\quad\text{and in exactly the same way}\quad[bL]=0,\ [cL]=0\ \text{etc.}
\]
Further, multiplying equations (2) in turn by \(\frac Lp,\ \frac{L'}{p'},\ \frac{L''}{p''}\) etc. and adding, we find
\[
\frac{lL}{p}+\frac{l'L'}{p'}+\frac{l''L''}{p''}+\text{etc.}=\frac{LL}{p}+\frac{L'L'}{p'}+\frac{L''L''}{p''}+\text{etc.}
\]
from which we obtain a \textit{second} expression for the weight,
\[
P=\frac{1}{\dfrac{lL}{p}+\dfrac{l'L'}{p'}+\dfrac{l''L''}{p''}+\text{etc.}}
\]
Finally, multiplying equations (2) in turn by \(\frac lp,\ \frac{l'}{p'},\ \frac{l''}{p''}\) etc. and adding, we arrive at a \textit{third} expression for the weight
\[
P=\frac{1}{[al]x^0+[bl]y^0+[cl]z^0+\text{etc.}+[ll]}
\]
if, in keeping with the other notation, we set
\[
\frac{ll}{p}+\frac{l'l'}{p'}+\frac{l''l''}{p''}+\text{etc.}=[ll]
\]
From this, with the aid of equations (1), it is easy to pass to a \textit{fourth expression}, which we give as follows:
\[
\begin{aligned}
\frac1P=[ll]&-[aa]x^0x^0-[bb]y^0y^0-[cc]z^0z^0-\text{etc.}\\
&-2[ab]x^0y^0-2[ac]x^0z^0-2[bc]y^0z^0-\text{etc.}
\end{aligned}
\]
\SourceLeaf{070}
\Article{8}
The general solution that we have developed so far is chiefly suited to the case in which \textit{one} unknown quantity depending on observed quantities is to be determined. But whenever the most plausible values of several unknown quantities depending on the same observations are required, or whenever it is still uncertain which particular unknown quantities ought to be derived from the observations, it will be appropriate to prepare the observations in another way, which we now proceed to develop.

We shall regard the quantities \(x,y,z\) etc. as indeterminates, set
\[
\left.\begin{aligned}\relax
[aa]x+[ab]y+[ac]z+\text{etc.}&=\xi\\
[ab]x+[bb]y+[bc]z+\text{etc.}&=\eta\\
[ac]x+[bc]y+[cc]z+\text{etc.}&=\zeta
\end{aligned}\right\}\ (6)
\]
etc., and suppose that elimination gives
\[
\left.\begin{aligned}\relax
[\alpha\alpha]\xi+[\alpha\beta]\eta+[\alpha\gamma]\zeta+\text{etc.}&=x\\
[\beta\alpha]\xi+[\beta\beta]\eta+[\beta\gamma]\zeta+\text{etc.}&=y\\
[\gamma\alpha]\xi+[\gamma\beta]\eta+[\gamma\gamma]\zeta+\text{etc.}&=z
\end{aligned}\right\}\ (7)
\]
etc.

Here it must first of all be observed that coefficients placed symmetrically must be equal, for example
\[
\begin{aligned}\relax
[\beta\alpha]&=[\alpha\beta]\\
[\gamma\alpha]&=[\alpha\gamma]\\
[\gamma\beta]&=[\beta\gamma]\\
\text{etc.}&
\end{aligned}
\]
This follows at once from the general theory of elimination in linear equations, but it will also be proved directly below without that theory.

We shall therefore have
\[
\left.\begin{aligned}
x^0&=-[\alpha\alpha].[al]-[\alpha\beta].[bl]-[\alpha\gamma].[cl]-\text{etc.}\\
y^0&=-[\alpha\beta].[al]-[\beta\beta].[bl]-[\beta\gamma].[cl]-\text{etc.}\\
z^0&=-[\alpha\gamma].[al]-[\beta\gamma].[bl]-[\gamma\gamma].[cl]-\text{etc.}
\end{aligned}\right\}\ (8)
\]
\begin{center}etc.\end{center}
from which, if we set
\SourceLeaf{071}
\[
\left.\begin{aligned}\relax
[\alpha\alpha]\mathfrak A+[\alpha\beta]\mathfrak B+[\alpha\gamma]\mathfrak C+\text{etc.}&=A\\
[\alpha\beta]\mathfrak A+[\beta\beta]\mathfrak B+[\beta\gamma]\mathfrak C+\text{etc.}&=B\\
[\alpha\gamma]\mathfrak A+[\beta\gamma]\mathfrak B+[\gamma\gamma]\mathfrak C+\text{etc.}&=C
\end{aligned}\right\}\ (9)
\]
etc., we obtain
\[
K=k-A[al]-B[bl]-C[cl]-\text{etc.}
\]
or if we further set
\[
\left.\begin{aligned}
aA+bB+cC+\text{etc.}&=p\varepsilon\\
a'A+b'B+c'C+\text{etc.}&=p'\varepsilon'\\
a''A+b''B+c''C+\text{etc.}&=p''\varepsilon''
\end{aligned}\right\}\ (10)
\]
etc., we shall have
\[
K=k-l\varepsilon-l'\varepsilon'-l''\varepsilon''-\text{etc.}\qquad(11)
\]
\clearpage
